\section{Repetidores y regeneración}

Supongamos se transmite una onda digital de banda base por un
cable largo. Al propagarse la onda se introduce atenuación, que
debe compensarse por amplificación posterior. Una estrategia
comúnmente usada es dividir el cable en tramos y amplificar en
cada tramo. Ejemplo de uso: cajas subterráneas en telefonía.

\begin{figure}[htb]
 \begin{center}\begin{minipage}{\textwidth}
  \bpgf{Imagenes/repReg/repReg-1}
  \begin{blockdiagram}
    \matrix[row sep=5mm, column sep=10mm]{
        & \node[etiqueta] (ruido1) {$n_1$}; & & \node[etiqueta] (ruido2) {$n_2$}; & & & \node[etiqueta] (ruidom) {$n_m$};\\
        \node[etiqueta] (inicio) {$x(t)$}; & \node[operator] (sum1) {$+$}; & \node[block] (amp1) {AMP}  node [above right=5mm] {$y_1(t)$} node [below=3mm] {\begin{minipage}{1.5cm}\centering ganancia\\ $\frac{1}{a}$\end{minipage}}; & \node[operator] (sum2) {$+$}; & \node[block] (amp2) {AMP}; & \node[etiqueta] (puntos) {$\ldots$}; & \node[operator] (summ) {$+$}; & \node[block] (ampm) {AMP}; & \node[etiqueta] (salida) {$y(t)$};\\
    };
    \draw [->,ultra thick] (inicio) -- node [anchor=north] {\begin{minipage}{20mm}\centering TR.1\\ atenuación $a$ ($a\ll 1$)\end{minipage}} (sum1);
    \draw [->] (sum1) -- (amp1);
    \draw [->] (ruido1) -- (sum1);
    \draw [->,ultra thick] (amp1) -- node [anchor=north] {TR. 2} (sum2);
    \draw [->] (ruido2) -- (sum2);
    \draw [->] (sum2) -- (amp2);
    \draw [->,ultra thick] (amp2) -- (puntos);
    \draw [->,ultra thick] (puntos) -- (summ);
    \draw [->] (ruidom) -- (summ);
    \draw [->] (summ) -- (ampm);
    \draw [->] (ampm) -- (salida);
  \end{blockdiagram}
  \endpgfgraphicnamed\end{minipage}
 \end{center}

\end{figure}
\FloatBarrier

Suponemos que cada etapa tiene ruido aleatorio de varianza
$E[n_i^2] = \sigma^2$.

\[y_1 = \frac{1}{a}(a x + n_1) = x+ \frac{n_1}{a} \qquad \Longrightarrow \left(\frac{S}{N}\right)_1 =
\frac{\langle x^2 \rangle}{\frac{\langle n_1^2\rangle}{a^2}} = a^2
\frac{\langle x^2\rangle}{\sigma^2}.\]

Para la segunda etapa:
\[y_2 = \frac{1}{a}(a y_1 +n_2) = y_1 + \frac{n_2}{a} = x + \frac{n_1+n_2}{a}.\]

Por el mismo razonamiento,
\[y(t) = y_m(t) = x(t) + \frac{n_1+n_2+\cdots+n_m}{a}.\]

\begin{align*}
 \left(\frac{S}{N}\right)_m & = \frac{\langle x^2 \rangle}{\frac{1}{a^2} E\left[(n_1+n_2+\cdots+n_m)^2\right]} =
  \frac{a^2 \langle x^2\rangle}{m \sigma^2}, \qquad \text{suponiendo ruidos independientes.}
%\Rightarrow                   & =
\end{align*}
\[{\col{ \left(\frac{S}{N}\right)_m = \frac{1}{m}
\left(\frac{S}{N}\right)_1}}\]

 Vemos que la relación señal-ruido se degrada linealmente en $m$. Por lo tanto
la $P_e$ se degrada fuertemente (por ejemplo,
$P_e=Q\left(\sqrt{\frac{S}{N}}\right)$ para el caso polar, con
pulsos ideales).

\begin{ejemplo}[numérico] Señal polar, $\left(\frac{S}{N}\right)_1 = 100 \  (20\,dB)$, $m=20$
repetidores.
  \[ \Rightarrow \left(\frac{S}{N}\right)_m = 5 \Rightarrow P_e  = Q(\sqrt{5}) = 1.29\cdot 10^{-2}, \quad\text{muy alto.}\]
 \end{ejemplo}

\subsection*{Repetidor regenerativo}

En vez de solo amplificar (señal y ruido), realiza una detección y
regenera la señal PAM. Sea $p$ la probabilidad de error de una
etapa, $p=Q\left(\sqrt{\left(\frac{S}{N}\right)_1}\right)$.

Suponemos $p\ll 1$ y errores independientes en cada etapa.
Calculamos la probabilidad de error de $m$ etapas. Notar que
\[1 - P_e(m\text{ etapas})  = (1-P_e(1 \text{ etapa}))^m \]
(es decir, para no tener errores en un símbolo exijo no tenerlos
en ninguna de las $m$ detecciones independientes). Por tanto
\[P_e(m\text{ etapas}) = 1 -(1-p)^m  \approx 1-(1-mp) = mp.\]
En este caso, $P_e$ se deteriora linealmente con $m$.

Volviendo al ejemplo, $p=Q(10) = 7.7\cdot 10^{-24}$ para una
etapa, por tanto
\[ P_e (20 \text{ etapas}) \approx 20 p = 1.5\cdot 10^{-22},\]
muchísimo más bajo que en el caso no regenerativo.

\paragraph{Moraleja}$\phantom{1}$
\begin{itemize}
 \item Es mejor regenerar.
 \item Esta es una de las mayores ventajas de la comunicación digital por sobre la analógica.
\end{itemize}
